\documentclass[11pt,a4paper]{article}

\usepackage[utf8]{inputenc}
\usepackage[T1]{fontenc}
\usepackage[polish]{babel}
\usepackage{lmodern}
\usepackage{microtype}
\usepackage[a4paper,margin=2cm]{geometry}
\usepackage{array}
\usepackage{booktabs}
\usepackage{tabularx}
\usepackage{enumitem}
\usepackage{xcolor}
\usepackage[colorlinks=true,linkcolor=blue!50!black,urlcolor=blue!50!black]{hyperref}

\newcolumntype{L}{>{\raggedright\arraybackslash}X}
\newcolumntype{P}[1]{>{\raggedright\arraybackslash}p{#1}}
\setlist{itemsep=2pt,topsep=4pt}
\setlength{\parindent}{0pt}
\setlength{\parskip}{6pt}
\renewcommand{\arraystretch}{1.25}

\title{Założenia projektowe\\[4pt]\large Sklep internetowy z elektroniką}
\author{}
\date{}

\begin{document}

\maketitle
\vspace{-2em}
\tableofcontents

\section{Cel i zakres}

Celem jest serwis WWW -- sklep internetowy z elektroniką użytkową. Oprócz zwykłej sprzedaży będzie on dobierał
klientowi produkty na podstawie jego zachowania i pomagał mu wybrać sprzęt w rozmowie z asystentem AI.

Poza zakresem: prawdziwe płatności online (płatność będzie symulowana), integracja z firmami kurierskimi,
wysyłka e-maili, opinie o produktach.

\section{Technologie}

{\small
\begin{tabularx}{\textwidth}{P{2.8cm}LL}
\toprule
\textbf{Element} & \textbf{Technologia} & \textbf{Zastosowanie} \\
\midrule
Frontend & React, TypeScript, Mantine, Tailwind CSS & strona sklepu i panel administratora \\
Backend & Java, Spring Boot & logika sklepu i inteligentnych usług \\
Baza danych & PostgreSQL z rozszerzeniem pgvector & dane sklepu i wektory opisów produktów \\
Logowanie & Keycloak (OAuth 2.0 / OpenID Connect) & rejestracja, logowanie, role \\
Model językowy & Claude Opus 5.5 (API Anthropic) & asystent AI \\
Wektory opisów & wielojęzyczny model bge-m3 uruchamiany lokalnie w Ollama & wyszukiwanie produktów po
znaczeniu opisu \\
Uruchamianie & Docker Compose & całe środowisko uruchamiane jedną komendą \\
\bottomrule
\end{tabularx}\par\medskip
}

\section{Role użytkowników}

{\small
\begin{tabularx}{\textwidth}{P{2.8cm}L}
\toprule
\textbf{Rola} & \textbf{Uprawnienia} \\
\midrule
Gość & przeglądanie i wyszukiwanie produktów, koszyk, bestsellery, czat z asystentem \\
Klient & jak gość + składanie zamówień, historia zamówień, strona konta, rekomendacje personalizowane \\
Administrator & zarządzanie produktami i stanami magazynowymi, zmiana statusów zamówień \\
\bottomrule
\end{tabularx}\par\medskip
}

\section{Dane produktów}

Każdy produkt będzie miał:
\begin{itemize}
  \item nazwę, markę, kategorię, cenę brutto i stan magazynowy,
  \item opis i zdjęcia,
  \item specyfikację do wyświetlenia na karcie produktu,
  \item \textbf{atrybuty techniczne}, z których korzystają dopasowanie (\ref{sec:dopasowanie}) i zestawy
        (\ref{sec:zestawy}),
  \item \textbf{ocenę jakości} w skali 1--100, według której będzie wybierany najlepszy zestaw.
\end{itemize}

Przykładowe grupy kategorii: audio (słuchawki, głośniki), smartfony i smartwatche, komputery (laptopy,
tablety), monitory i peryferia (myszy, klawiatury), akcesoria (etui, ładowarki, kable, stacje dokujące).

\filbreak
Zakładane atrybuty techniczne:

{\small
\begin{tabularx}{\textwidth}{P{3.6cm}L}
\toprule
\textbf{Kategoria} & \textbf{Atrybuty} \\
\midrule
smartfon & model, złącze ładowania \\
etui & model telefonu, do którego pasuje \\
ładowarka, kabel & złącze, moc \\
laptop & wyjścia wideo, USB-C z obrazem \\
monitor & wejścia wideo, przekątna, częstotliwość odświeżania \\
stacja dokująca & wymagane USB-C z obrazem \\
\bottomrule
\end{tabularx}\par\medskip
}

Sklep będzie miał przygotowane dane startowe, wystarczające do pokazania wszystkich funkcji.

\section{Sklep -- wymagania funkcjonalne}

{\small
\begin{tabularx}{\textwidth}{LP{1.5cm}}
\toprule
\textbf{Wymaganie} & \textbf{Rola} \\
\midrule
Rejestracja i logowanie przez Keycloak z własnym motywem strony logowania & gość \\
Katalog produktów ze stronicowaniem & wszyscy \\
Filtry (kategoria, marka, cena, łączność) i sortowanie (trafność, cena, nowość, popularność) & wszyscy \\
Wyszukiwarka działająca jak w punkcie~\ref{sec:wyszukiwanie} & wszyscy \\
Karta produktu: zdjęcia, cena, dostępność, opis, specyfikacja, rekomendacje
z punktu~\ref{sec:rekomendacje} & wszyscy \\
Koszyk: dodawanie, zmiana ilości, usuwanie, podsumowanie & wszyscy \\
Zamówienie: dane kontaktowe, adres, metoda dostawy (kurier, odbiór osobisty), płatność (BLIK, karta),
faktura & klient \\
Cena zamówienia liczona po stronie serwera z aktualnych cen; rezerwacja towaru przy zamówieniu, zwolnienie
przy nieudanej płatności lub anulowaniu & klient \\
Historia zamówień i szczegóły zamówienia ze statusem & klient \\
Strona konta: dane klienta, ostatnie zamówienia, podsumowanie zakupów, zapisany adres & klient \\
Panel: dodawanie, edycja i usuwanie produktów, zmiana stanów magazynowych & admin \\
Panel: lista zamówień, zmiana statusu & admin \\
\bottomrule
\end{tabularx}\par\medskip
}

\filbreak
Statusy zamówienia i dozwolone przejścia:

{\small
\begin{tabularx}{\textwidth}{P{4cm}LP{3.2cm}}
\toprule
\textbf{Status} & \textbf{Może przejść do} & \textbf{Kto zmienia} \\
\midrule
Oczekuje na płatność & Opłacone, Płatność nieudana & system płatności \\
Opłacone & Wysłane, Anulowane & administrator \\
Wysłane & -- & -- \\
Płatność nieudana & -- & -- \\
Anulowane & -- & -- \\
\bottomrule
\end{tabularx}\par\medskip
}

\section{Inteligentne usługi}

\subsection{Rejestrowanie aktywności}\label{sec:aktywnosc}

Sklep będzie zapisywał zdarzenia zalogowanych klientów: wyświetlenie produktu, wyszukiwanie (z treścią
zapytania), dodanie do koszyka i zakup. Na tych danych będą się opierać rekomendacje i ich ocena.

\subsection{Symulator klientów}\label{sec:symulator}

Sklep nie będzie miał prawdziwych klientów, więc dane o aktywności wygeneruje symulator. Przy każdym
uruchomieniu da te same wyniki, żeby ocena jakości była powtarzalna.

Każdy symulowany klient będzie miał jeden typ:

{\small
\begin{tabularx}{\textwidth}{P{3.6cm}L}
\toprule
\textbf{Typ klienta} & \textbf{Czym się interesuje} \\
\midrule
gracz & laptopy do gier, monitory, myszy, klawiatury, słuchawki \\
biegacz & słuchawki sportowe, smartwatche \\
student & laptopy, tablety, słuchawki, myszy \\
fan muzyki & słuchawki, głośniki \\
pracownik biurowy & laptopy, monitory, stacje dokujące, klawiatury, myszy \\
fan smartfonów & smartfony, smartwatche, akcesoria do telefonów \\
\bottomrule
\end{tabularx}\par\medskip
}

Zasady zachowania symulowanych klientów:
\begin{itemize}
  \item większość oglądanych produktów pochodzi z kategorii ich typu, reszta jest losowa,
  \item każdy klient ma ulubione marki, które wybiera częściej,
  \item każdy klient ma budżet zależny od typu,
  \item tylko część obejrzanych produktów trafia do koszyka i tylko część koszyków kończy się zakupem,
  \item po zakupie smartfona lub laptopa klient często dokupuje pasujące akcesorium
        (punkt~\ref{sec:dopasowanie}).
\end{itemize}

\subsection{Rekomendacje produktów}\label{sec:rekomendacje}

{\footnotesize
\begin{tabularx}{\textwidth}{P{2.7cm}P{2.6cm}LP{3.6cm}}
\toprule
\textbf{Sekcja} & \textbf{Gdzie} & \textbf{Jak liczona} & \textbf{Gdy brak danych} \\
\midrule
„Polecane dla Ciebie” & strona główna zalogowanego & filtrowanie kolaboratywne (opis niżej) & mało
historii: reguła kategoria/marka; brak historii: bestsellery \\
„Na podstawie oglądanych” & strona główna zalogowanego & produkty z kategorii ostatnio oglądanych,
w podobnej cenie & sekcja ukryta \\
„Bestsellery” & strona główna gościa & najczęściej kupowane produkty z ostatniego miesiąca & najnowsze
produkty \\
„Często kupowane razem” & karta produktu & produkty, które najczęściej trafiały do tych samych zamówień &
pasujące akcesoria z punktu~\ref{sec:dopasowanie} \\
„Podobne produkty” & karta produktu & produkty z tej samej kategorii o najbardziej podobnym opisie
(\ref{sec:wyszukiwanie}) i zbliżonej cenie & ta sama kategoria, najbliższa cena \\
\bottomrule
\end{tabularx}\par\medskip
}

Rekomendacje nie będą pokazywać produktów niedostępnych ani już kupionych przez klienta.

\textbf{Reguła kategoria/marka} -- prosty punkt odniesienia przy ocenie jakości. Polecane są produkty
z kategorii i marek, które klient już kupował, przy czym zgodność kategorii liczy się bardziej niż zgodność
marki.

\textbf{Filtrowanie kolaboratywne} -- „klienci, którzy interesowali się tym co Ty, interesowali się
też\dots”:
\begin{enumerate}
  \item Każdy kontakt klienta z produktem ma wagę: zakup waży więcej niż dodanie do koszyka, a to więcej niż
        wyświetlenie.
  \item Dwa produkty są podobne, gdy interesują się nimi ci sami klienci (podobieństwo cosinusowe).
  \item Klient dostaje produkty najbardziej podobne do tych, z którymi miał kontakt, z uwzględnieniem wag
        kontaktów.
  \item Podobieństwa będą przeliczane okresowo, a nie przy każdym wyświetleniu strony.
\end{enumerate}

\subsection{Wyszukiwanie po opisie potrzeb}\label{sec:wyszukiwanie}

Wyszukiwarka połączy dwa sposoby szukania:

{\small
\begin{tabularx}{\textwidth}{P{2.4cm}LP{4.2cm}}
\toprule
\textbf{Sposób} & \textbf{Na czym polega} & \textbf{Przykład, który ma zadziałać} \\
\midrule
po słowach & podobieństwo zapytania do nazwy, marki i kategorii, odporne na literówki & „sony xm5” \\
po znaczeniu & opis produktu i zapytanie zamieniane na wektory, porównanie ich podobieństwa & „cicha
klawiatura do biura” \\
\bottomrule
\end{tabularx}\par\medskip
}

\begin{itemize}
  \item Wyniki obu sposobów łączymy metodą RRF: na górę trafiają produkty, które są wysoko na obu listach.
  \item Kwoty w zapytaniu („do 300 zł”, „od 1000 do 2000 zł”) działają jak filtr ceny.
  \item Filtry i sortowanie z katalogu działają też na wynikach wyszukiwania.
\end{itemize}

\subsection{Dopasowanie akcesoriów}\label{sec:dopasowanie}

Sklep sprawdzi, czy akcesoria pasują do urządzenia, według stałych zasad opartych na atrybutach
z rozdziału 4:

{\small
\begin{tabularx}{\textwidth}{P{5cm}L}
\toprule
\textbf{Para} & \textbf{Zasada} \\
\midrule
etui -- smartfon & etui pasuje do modelu smartfona \\
ładowarka, kabel -- smartfon & to samo złącze \\
stacja dokująca -- laptop & laptop ma USB-C z obrazem \\
monitor -- laptop & mają wspólne złącze wideo \\
\bottomrule
\end{tabularx}\par\medskip
}

Dopasowanie będzie widoczne w trzech miejscach:
\begin{itemize}
  \item karta smartfona i laptopa -- sekcja „Pasujące akcesoria” z akcesoriami, które pasują do tego
        urządzenia,
  \item koszyk -- ostrzeżenie z powodem, gdy akcesorium nie pasuje do urządzenia w koszyku, np. „Etui pasuje
        do iPhone 15, a w koszyku jest iPhone 16”,
  \item czat z asystentem -- asystent sprawdzi dopasowanie, zanim zaproponuje akcesoria lub zestaw.
\end{itemize}

\subsection{Zestawy w budżecie}\label{sec:zestawy}

Zestawy będą dostępne w czacie z asystentem. Klient napisze, czego potrzebuje i ile chce wydać, np.
\emph{„mam 5000 zł, potrzebuję stanowiska do pracy z laptopem”}. Asystent pokaże zestaw jako karty produktów
z łączną ceną i przyciskiem „Dodaj cały zestaw”.

Zestaw będzie składany według jednego z szablonów:

{\small
\begin{tabularx}{\textwidth}{P{4.2cm}L}
\toprule
\textbf{Szablon} & \textbf{Części (od najważniejszej)} \\
\midrule
stanowisko z laptopem & laptop, monitor, stacja dokująca, mysz, klawiatura \\
smartfon z akcesoriami & smartfon, etui, ładowarka, słuchawki \\
\bottomrule
\end{tabularx}\par\medskip
}

Sposób wyboru:
\begin{enumerate}
  \item Dla każdej części wybieramy kilku kandydatów z różnych przedziałów cenowych, najlepszych pod względem
        oceny jakości i zgodnych z wymaganiami klienta (np. „wolę Samsunga”).
  \item Sprawdzamy kombinacje kandydatów i odrzucamy te, które przekraczają budżet lub nie pasują do siebie
        (\ref{sec:dopasowanie}).
  \item Wybieramy kombinację o najwyższej łącznej ocenie jakości, w której ważniejsze części liczą się
        bardziej.
  \item Gdy nic nie mieści się w budżecie, asystent pokaże najtańszy pasujący zestaw i brakującą kwotę.
\end{enumerate}

\subsection{Asystent AI}\label{sec:asystent}

W rogu każdej strony będzie przycisk czatu z asystentem, a na telefonie czat otworzy się na pełnym ekranie.
Odpowiedź będzie się pojawiać na bieżąco, w trakcie pisania. Proponowane produkty wyświetlą się jako karty
z ceną i przyciskiem „Dodaj do koszyka”, a zestaw -- z łączną ceną i przyciskiem „Dodaj cały zestaw”.

Asystent ma:
\begin{itemize}
  \item odpowiadać po polsku, krótko, z kilkoma propozycjami naraz i jednozdaniowym uzasadnieniem każdej,
  \item dopytać o budżet lub zastosowanie, gdy ich nie zna,
  \item brać nazwy, ceny i dostępność wyłącznie z danych sklepu, nigdy z własnej wiedzy,
  \item poprawiać propozycję na prośbę klienta, np. „a coś tańszego?”, „wolę Lenovo”,
  \item porównywać produkty, gdy klient o to poprosi,
  \item uwzględniać ostatnie zakupy zalogowanego klienta, np. nie proponować drugiej myszy,
  \item odmawiać rozmowy niezwiązanej z zakupami elektroniki,
  \item nie zmieniać cen, nie dawać rabatów i nie ujawniać swoich instrukcji.
\end{itemize}

Model językowy będzie rozumiał wiadomość klienta i pisał odpowiedź, a dane będzie
brał z funkcji sklepu:

{\small
\begin{tabularx}{\textwidth}{P{6cm}L}
\toprule
\textbf{Funkcja dostępna dla asystenta} & \textbf{Na czym się opiera} \\
\midrule
wyszukanie produktów & wyszukiwarka z punktu~\ref{sec:wyszukiwanie} \\
szczegóły produktu & aktualna cena, dostępność, specyfikacja \\
sprawdzenie dopasowania & zasady z punktu~\ref{sec:dopasowanie} \\
złożenie zestawu w budżecie & algorytm z punktu~\ref{sec:zestawy} \\
ostatnie zakupy klienta & historia zamówień zalogowanego klienta \\
pokazanie produktów jako kart & wyświetlenie kart w oknie czatu \\
\bottomrule
\end{tabularx}\par\medskip
}

Ceny i zestawy będzie liczył sklep, a nie model, dlatego asystent nie poda nieistniejącego produktu ani złej
ceny. Asystent sam niczego nie doda do koszyka -- zrobi to klient przyciskiem na karcie. Długość wiadomości
i liczba wiadomości w rozmowie będą ograniczone.

\subsection{Ocena jakości}\label{sec:ocena}

Każda inteligentna usługa zostanie zmierzona i porównana z prostszym rozwiązaniem. Wyniki trafią do
sprawozdania.

Rekomendacje i wyszukiwanie będą oceniane trzema miarami liczonymi dla 10 pierwszych pozycji listy:

{\small
\begin{tabularx}{\textwidth}{P{3cm}L}
\toprule
\textbf{Miara} & \textbf{Pytanie} \\
\midrule
precision@10 & jaka część z 10 pokazanych produktów to trafienia? \\
recall@10 & jaką część wszystkich trafnych produktów udało się pokazać? \\
NDCG@10 & czy trafienia są na górze listy? (1 = idealna kolejność, 0 = brak trafień) \\
\bottomrule
\end{tabularx}\par\medskip
}

\section{Bezpieczeństwo}

\begin{itemize}
  \item \textbf{Logowanie} -- wyłącznie przez Keycloak (OAuth 2.0 / OpenID Connect). Aplikacja nie będzie
        przechowywać haseł.
  \item \textbf{Uprawnienia} -- backend będzie sprawdzał tożsamość i rolę przy każdym żądaniu, zgodnie
        z tabelą ról z rozdziału 3.
  \item \textbf{Izolacja danych} -- klient zobaczy tylko własne zamówienia.
  \item \textbf{Dane wejściowe} -- każde żądanie będzie walidowane, a baza chroniona przed SQL injection.
  \item \textbf{Ceny} -- liczone wyłącznie na serwerze, niezależnie od tego, co wyśle przeglądarka.
  \item \textbf{Asystent AI} -- będzie tylko odczytywał dane sklepu, nie zobaczy danych osobowych klientów
        i nie będzie mógł zmieniać cen, koszyka ani zamówień. Klucz do API modelu językowego będzie tylko
        na serwerze.
  \item \textbf{Testy} -- testy automatyczne sprawdzą dostęp do każdej funkcji bez logowania, jako klient
        i jako administrator.
\end{itemize}

\section{Wymagania niefunkcjonalne}

\begin{itemize}
  \item Typowa odpowiedź serwera poniżej 0,5 s.
  \item Pierwsze słowa odpowiedzi asystenta w ciągu kilku sekund.
  \item Interfejs responsywny (telefon, tablet, komputer), w języku polskim.
  \item Całe środowisko uruchamiane jedną komendą na Windows, Linux i macOS.
  \item Opis API generowany automatycznie z kodu backendu, a na jego podstawie klient API we frontendzie.
  \item Testy integracyjne backendu na prawdziwej bazie PostgreSQL.
\end{itemize}

\section{Podział pracy}

Każda osoba odpowiada za swój obszar w całości: backend, frontend, testy i część sprawozdania. Każda ma
porównywalną ilość pracy przy sklepie i przy inteligentnych usługach.

{\footnotesize
\begin{tabularx}{\textwidth}{P{1cm}P{2.4cm}LL}
\toprule
\textbf{Osoba} & \textbf{Obszar} & \textbf{Sklep} & \textbf{Inteligentne usługi} \\
\midrule
A & zakupy i asystent & środowisko uruchomieniowe, koszyk, zamówienie z płatnością symulowaną & asystent AI
(\ref{sec:asystent}), ocena asystenta \\
B & katalog i wyszukiwanie & dane produktów z atrybutami, katalog, filtry, sortowanie, karta produktu &
wyszukiwanie (\ref{sec:wyszukiwanie}), dopasowanie akcesoriów (\ref{sec:dopasowanie}), zestawy
(\ref{sec:zestawy}), ocena wyszukiwania i zestawów \\
C & klienci i personalizacja & Keycloak (role, motyw logowania), historia zamówień, strona konta, panel
administratora & aktywność (\ref{sec:aktywnosc}), symulator (\ref{sec:symulator}), rekomendacje
(\ref{sec:rekomendacje}), ocena rekomendacji \\
\bottomrule
\end{tabularx}\par\medskip
}

Funkcje, z których korzysta asystent, dostarczają osoby B i C. Osoba A podłącza je do asystenta.

\section{Kolejność prac}

{\footnotesize
\begin{tabularx}{\textwidth}{P{1.4cm}LLL}
\toprule
\textbf{Tydzień} & \textbf{Osoba A} & \textbf{Osoba B} & \textbf{Osoba C} \\
\midrule
1--2 & środowisko uruchomieniowe, szkielet backendu i frontendu & model danych produktu z atrybutami,
pierwsze dane & Keycloak, role, logowanie we frontendzie \\
3--5 & koszyk, zamówienie, płatność & katalog, filtry, karta produktu, komplet danych & strona konta, panel
administratora (produkty), rejestrowanie aktywności \\
6--8 & czat z asystentem: rozmowa, odpowiedź na bieżąco, wyszukiwanie produktów & wektory opisów,
wyszukiwanie po słowach i po znaczeniu, łączenie wyników & historia zamówień, panel administratora
(zamówienia), symulator klientów \\
9--11 & pozostałe funkcje asystenta, karty produktów i zestawów w czacie & dopasowanie akcesoriów, zestawy
w budżecie, „Pasujące akcesoria” & bestsellery, reguła kategoria/marka, filtrowanie kolaboratywne,
pozostałe sekcje rekomendacji \\
12--13 & rozmowy testowe, ocena asystenta, poprawki & zapytania testowe, ocena wyszukiwania i zestawów &
ocena rekomendacji \\
14--15 & sprawozdanie, scenariusze pokazowe, prezentacja & sprawozdanie, scenariusze pokazowe, prezentacja &
sprawozdanie, scenariusze pokazowe, prezentacja \\
\bottomrule
\end{tabularx}\par\medskip
}

\end{document}
